% !TEX root = ../thesis.tex

\chapter{Conclusion, Discussion, and Future Work}
\label{chp:conclusion}

\section{Summary of Contributions}

This project adapted the RoboDuet decentralised multi-agent RL framework
from the Unitree Go1+Z1 platform to the heavier and more capable Unitree
B2+Z1 system.
The main technical contributions are:

\begin{enumerate}
  \item \textbf{Platform migration}: complete adaptation of URDF loading,
        DOF indexing, PD gain calibration, and simulation configuration for
        B2Z1.  The 19-DOF system structure---12 leg, 6 arm, 1 gripper---was
        fully characterised, including the discovery that the gripper
        (DOF~18) requires direct state injection due to zero PD gains.
  \item \textbf{Forward reach validation}: identification that the arm
        policy's forward reach direction is encoded by the spherical command
        $(l, p, \psi) = (0.55, 0.25, 0)$, yielding EE position
        $(0.533, 0, 0.516)$\,m, in contrast to scripted joint angle
        approaches that produce backward reach on B2Z1.
  \item \textbf{Three reproducible demonstrations}: fixed-base arm
        manipulation (10\,s), standing quadruped with arm reaching (10\,s),
        and a full fixed-base grasp-and-lift sequence (20\,s, 30\,cm lift),
        all produced as H.264 MP4 videos at 1280$\times$720 resolution.
  \item \textbf{Stability characterisation}: base height variation of only
        $\pm$3\,cm in the free-standing demonstration; kinematically fixed
        at 0.340\,m in fixed-base demonstrations, confirming that the
        decentralised two-policy architecture is stable on B2Z1 without
        retraining.
\end{enumerate}

\section{Limitations}

\textbf{Simulation only.}  All results are in the IsaacGym simulator.
Sim-to-real transfer for the heavier B2 body has not been validated;
the B2's different foot geometry, motor characteristics, and higher inertia
may require additional domain randomisation or system-identification steps.

\textbf{No physical grasping.}  The grasp-and-lift demonstration uses
kinematic box teleportation rather than rigid-body contact.  The gripper
does not physically close around the box; the visual effect is produced by
a hybrid strategy that tracks the live gripper XY position while
independently animating the box Z coordinate upward by 30\,cm.
True physical grasping would require a contact-rich simulation or
force/torque feedback.

\textbf{Standing only.}  The demonstration scripts fix locomotion velocity
commands to zero.  Concurrent walking and manipulation---the primary
use-case of a loco-manipulator---has not been demonstrated on B2Z1.
The dog policy can walk with arm policy active, but this requires additional
tuning of the camera angle and scripting of a walking-and-reaching task.

\textbf{Single checkpoint.}  Only one pre-trained checkpoint
(\texttt{ac\_weights\_last\_arm.pt}) was evaluated.  A systematic evaluation
across multiple seeds or training iterations was outside the scope of this
term.

\section{Closing Remarks and Future Directions}

The results demonstrate that RoboDuet's decentralised policies and
spherical command space transfer effectively to the B2Z1 platform
despite its significantly higher body mass compared with Go1.
The engineering challenges encountered---backward reach diagnosis,
gripper DOF bypass, box teleport stability---provide practical insights
for deploying RL-based loco-manipulation on new hardware.

%% ========================================================================= %%
%%  Plans for Term 3 (Summer Session 2025/26)                                %%
%% ========================================================================= %%
\section*{Plans for the Second Term (Summer Session 2025/26)}
\addcontentsline{toc}{section}{Plans for the Second Term}

Building on the simulation results, the following milestones are planned
for the Summer Session:

\begin{enumerate}[nosep]
  \item \textbf{Sim-to-Real Transfer (Weeks 1--4):} deploy the arm
        policy on the physical B2Z1 via the Unitree SDK; perform system
        identification and validate EE reaching on hardware.
  \item \textbf{Contact-Rich Grasping (Weeks 3--6):} enable
        physics-based gripper--object contact in IsaacGym; add a grasp
        stability reward and retrain the arm policy.
  \item \textbf{Walking + Manipulation (Weeks 5--8):} develop a combined
        demo where the robot walks to an object, grasps it, and carries
        it to a target location.
  \item \textbf{Vision-Based Grasping (Weeks 7--10):} replace spherical
        commands with RGB-D observations for autonomous grasp planning.
  \item \textbf{B2Z1-Specific Retraining:} train from scratch with
        B2-specific domain randomisation to potentially outperform
        zero-shot transfer.
\end{enumerate}

\noindent
A final report documenting the sim-to-real results will be submitted at
the end of the Summer Session.
